% STARTHERE

La struttura seguente ricalca la guideline di reporting STROBE (i
numeri tra parentesi graffe, che si consiglia di lasciare, indicano
gli item della stessa); ai fini della compilazione delle parti di
propria competenza si consiglia di fare riferimento allo
\href{https://www.ncbi.nlm.nih.gov/pmc/articles/pmid/17941715/}{STROBE
  explanation and elaboration}

\section{Title \{1a\}}

{\em Indicate the study's design with a commonly used term in the title or
the abstract}

\section{Abstract \{1b\}}

{\em Provide in the abstract an informative and balanced summary of what
was done and what was found}

\section{Introduction}

\subsection{Background/rationale \{2\}}

{\em Explain the scientific background and rationale for the investigation
being reported}

\subsection{Objectives \{3\}}

{\em State specific objectives, including any prespecified hypotheses}

\section{Methods}

\subsection{Study design \{4\}}

{\em Present key elements of study design early in the paper}

\subsection{Setting \{5\}}

{\em Describe the setting, locations, and relevant dates, including periods
of recruitment, exposure, follow-up, and data collection}

\subsection{Participants}

\subsubsection{Eligibility criteria \{6a\}}

{\em Cohort study — Give the eligibility criteria, and the sources and
methods of selection of participants. Describe methods of follow-up.\\
Case-control study — Give the eligibility criteria, and the sources and
methods of case ascertainment and control selection. Give the
rationale for the choice of cases and controls.\\
Cross-sectional study — Give the eligibility criteria, and the sources
and methods of selection of participants}

\subsubsection{Matching \{6b\}}

{\em Cohort study — For matched studies, give matching criteria and number
of exposed and unexposed.\\
Case-control study — For matched studies, give matching criteria and
the number of controls per case}

\subsection{Variables \{7\}}

{\em Clearly define all outcomes, exposures, predictors, potential
confounders, and effect modifiers. Give diagnostic criteria, if
applicable }

\subsection{Data sources/measurement \{8\}}

{\em For each variable of interest, give sources of data and details of
methods of assessment (measurement). Describe comparability of
assessment methods if there is more than one group}

\subsection{Bias \{9\}}

{\em Describe any efforts to address potential sources of bias }

\subsection{Study size \{10\}}

{\em Explain how the study size was arrived at}

\subsection{Quantitative variables \{11\}}

{\em Explain how quantitative variables were handled in the analyses. If
applicable, describe which groupings were chosen and why}

\subsection{Statistical methods}

\subsubsection{Main methods \{12a\}}

{\em Describe all statistical methods, including those used to control for
confounding}

\subsubsection{Methods for subgroups/interactions \{12b\}}

{\em Describe any methods used to examine subgroups and interactions}

\subsubsection{Missing data \{12c\}}

{\em Explain how missing data were addressed}

\subsubsection{Methods taking account of sampling strategy \{12d\}}

{\em Cohort study - If applicable, explain how loss to follow-up was
addressed.\\
Case-control study - If applicable, explain how matching of cases and
controls was addressed.\\
Cross-sectional study — If applicable, describe analytical methods
taking account of sampling strategy
}

\subsubsection{Sensitivity analyses \{12e\}}

{\em Describe any sensitivity analyses}

\section{Results}

\subsection{Participants}

{\em In merito a questa sezione, give information separately for cases and
controls in case-control studies and, if applicable, for exposed and
unexposed groups in cohort and cross-sectional studies.}

\subsubsection{Number of individuals \{13a\}}

{\em Report numbers of individuals at each stage of study—eg numbers
potentially eligible, examined for eligibility, confirmed eligible,
included in the study, completing follow-up, and analysed}

\subsubsection{Reason for non-participation \{13b\}}

{\em Give reasons for non-participation at each stage}

\subsubsection{Flow diagram \{13c\}}

{\em Consider use of a flow diagram}

\subsection{Descriptive data}

{\em In merito a questa sezione, give information separately for cases and
controls in case-control studies and, if applicable, for exposed and
unexposed groups in cohort and cross-sectional studies.}

\subsubsection{Study participants\{14a\}}

{\em Give characteristics of study participants (eg demographic, clinical,
social) and information on exposures and potential confounders}

\subsubsection{Missing data \{14b\}}

{\em Indicate number of participants with missing data for each variable of
interest}

\subsubsection{Follow up (cohort studies) \{14c\}}

{\em Summarise follow-up time (eg, average and total amount)}

\subsection{Outcome data \{15\}}

{\em In merito a questa sezione, give information separately for cases and
controls in case-control studies and, if applicable, for exposed and
unexposed groups in cohort and cross-sectional studies.\\
Cohort study — Report numbers of outcome events or summary measures
over time.\\
Case-control study — Report numbers in each exposure category, or
summary measures of exposure.\\
Cross-sectional study — Report numbers of outcome events or summary
measures.
}

\subsection{Main results}

\subsubsection{Estimates \{16a\}}

{\em Give unadjusted estimates and, if applicable, confounder-adjusted
estimates and their precision (eg, 95\% confidence interval). Make
clear which confounders were adjusted for and why they were included}

\subsubsection{Variables categorization \{16b\}}

{\em Report category boundaries when continuous variables were categorized}

\subsubsection{Absolute risk \{16c\}}

{\em If relevant, consider translating estimates of relative risk into
absolute risk for a meaningful time period}

\subsection{Other analyses \{17\}}

{\em Report other analyses done—eg analyses of subgroups and interactions,
and sensitivity analyses}

\section{Discussion}

\subsection{Key results \{18\}}

{\em Summarise key results with reference to study objectives}

\subsection{Limitations \{19\}}

{\em Discuss limitations of the study, taking into account sources of
potential bias or imprecision. Discuss both direction and magnitude of
any potential bias}

\subsection{Interpretation \{20\}}

{\em Give a cautious overall interpretation of results considering
objectives, limitations, multiplicity of analyses, results from
similar studies, and other relevant evidence}

\subsection{Generalisability \{21\}}

{\em Discuss the generalisability (external validity) of the study results}

\section{Other information}

\subsection{Funding \{22\}}

{\em Give the source of funding and the role of the funders for the present
study and, if applicable, for the original study on which the present
article is based}


